\subsection{Kernel Suite}
\label{sec:meth:kernels}

Our benchmark suite covers five operator categories that capture the dominant computation patterns in modern deep learning: matrix multiplication, attention, convolution, normalization, and element-wise or reduction operators.
In total, the suite contains 22 kernels: three \textbf{GEMM} workloads (dense matmul, batched matmul, fused linear+activation), one \textbf{attention} workload (scaled dot-product / FlashAttention variants), one \textbf{convolution} workload (Conv2d, $1\times1$--$7\times7$ filters including depthwise and strided), two \textbf{normalization} workloads (LayerNorm and RMSNorm), and fifteen \textbf{element-wise or reduction} workloads.
Kernels are drawn from TritonBench~\cite{li2025tritonbench}, the TileLang example repository~\cite{wang2025tilelang}, and our own implementations for configurations absent from both sources.
Per-kernel provenance and full selection criteria are in \cref{app:repro:kernels}.


\subsection{Baseline Construction}
\label{sec:meth:baselines}

For each workload, we compare DSL implementations against the strongest available vendor-library path through PyTorch; per-category baseline specifications are in \cref{app:repro:baselines}.
For attention, we use PyTorch scaled dot-product attention with the FlashAttention backend; because the Triton and TileLang implementations are not algorithmically identical to this baseline, we report attention results separately.


\subsection{Profiling Setup}
\label{sec:meth:profiling}

We measure both end-to-end latency and hardware performance counters.
Latency drives all comparisons against library baselines; hardware counters support root-cause analysis (RQ2).

\paragraph{End-to-end throughput}
We measure host-synchronized GPU execution time; clocks are pinned to a fixed frequency throughout and run-to-run relative standard deviation is $\leq 0.9\%$ (measurement protocol in \cref{app:repro:meas}).

\paragraph{Hardware performance counters}
For root-cause analysis, we collect Nsight Compute~\cite{nsightcompute2024} profiles grounded in seven counters (definitions in \cref{app:repro:meas}), each from a single execution with stability verified within 2\% across five runs.

\paragraph{Correctness}
Before any timing, we validate every DSL kernel against its library baseline at per-dtype tolerances (\texttt{fp32}~$10^{-5}$, \texttt{fp16}~$10^{-3}$, \texttt{bf16}~$10^{-2}$, relaxed $2\times$ for reductions) and an edge-case suite (NaN/Inf/denormal/all-equal) across all 22 kernels with 0 crashes.
FP32 GEMM is excluded as silent TF32 lowering rather than a defect (\cref{sec:analysis:jit}).


\subsection{RQ3 Optimization Methodology}
\label{sec:meth:optimization}

To produce the optimized kernels evaluated in \cref{sec:mitigation}, we apply an agentic AI kernel-generation tool~\cite{ako2026} under a fixed protocol (correctness-gated iterations, large-input shapes, no language switching) to instantiate the root-cause-derived optimization patterns of \cref{sec:analysis}.
Optimized kernels are validated provenance-agnostically against the same per-dtype tolerances and Nsight Compute counters, and recovered performance is reported as a lower bound on user-space-recoverable gain.


\subsection{Metrics}
\label{sec:meth:metrics}

Our primary metric is \textbf{library efficiency},
\[
E_{\mathrm{lib}} = \frac{t_{\mathrm{lib}}}{t_{\mathrm{DSL}}} \times 100\%,
\]
where $t_{\mathrm{lib}}$ denotes the latency of the cuBLAS or cuDNN baseline and $t_{\mathrm{DSL}}$ denotes the latency of the corresponding DSL implementation under the same hardware and input configuration.
A value of 100\% indicates parity with the library baseline; lower values indicate that the DSL implementation is slower.
Secondary metrics (throughput, hardware efficiency, memory bandwidth utilization) are used qualitatively in root-cause analysis (RQ2) and are not tabulated separately.


\subsection{Evaluation-Gap and Heuristic Measurement}
\label{sec:meth:evalgap}

To test whether existing benchmarks admit slow kernels (RQ1), we run each naive DSL kernel through KernelBench's evaluator, which overrides the candidate's input generators with the reference's and checks output equivalence on randomized inputs; we record the pass/fail verdict and measured slowdown against the PyTorch baseline (harness details in \cref{app:repro:meas}).
To anchor the RQ3 heuristic in a baseline-independent signal, we model the bytes moved for memory-bound kernels or Floating-Point Operations Per Second (FLOPs) for compute-bound kernels, then divide the achieved work-rate by the relevant hardware peak (HBM bandwidth or FP16 Tensor-Core throughput) to obtain its roofline fraction.


\subsection{Experimental Setup}
\label{sec:meth:setup}

\paragraph{Hardware}
We measure on two primary architectures: the A100-SXM4-40GB (Ampere, \texttt{sm\_80}), used for suite magnitude (\cref{sec:evaluation}) and root-cause counters (\cref{sec:analysis}), and the GH200 (Grace Hopper, \texttt{sm\_90}), used for the evaluation-gap demonstration (\cref{sec:eval:passesslow}) and roofline heuristic (\cref{tab:roofline}).
Full  specifications and clock-lock settings are in \cref{app:repro:stack}.

\paragraph{Software}
We use PyTorch~2.8.0+cu128, Triton~3.4.0, TileLang~0.1.6.post1, bundled cuDNN, and Nsight Compute~2026.2.0; complete version list is in \cref{app:repro:stack}.
